\documentclass[10pt,a4paper]{amsart}
\usepackage[margin=19mm]{geometry}
\usepackage{amsmath,amssymb,mathtools}
\usepackage[scr=rsfs,cal=euler]{mathalfa}
\usepackage{xfrac}
\usepackage[unicode,hidelinks,bookmarks=false]{hyperref}
\newcommand{\ie}{i.e.\ }
\newcommand{\cf}{cf.\ }
\newcommand{\resp}{respectively\ }
\newcommand{\Subsection}{\subsection}
\providecommand{\nobreakdash}{\nobreak\hbox{-}}
\setlength{\emergencystretch}{2em}
\multlinegap0pt
\usepackage{siunitx}
\usepackage{siunitx}
\usepackage{siunitx}
\usepackage{siunitx}
\usepackage{amssymb}
\usepackage{mathtools}
\usepackage{siunitx}
\usepackage{dsfont}
\usepackage{xcolor}
\usepackage{tikz}
\usepackage{algorithm}\usepackage{algpseudocode}
\usepackage[shortlabels]{enumitem}
\newtheorem{lemma}{Lemma}
\newtheorem{assumption}{Assumption}
\usepackage{cleveref}
\crefname{lemma}{Lemma}{Lemmas}
\crefname{assumption}{Assumption}{Assumptions}
\newcommand{\R}{\mathbb{R}}
\DeclareMathOperator{\Sym}{sym}
\newcommand{\TenDeriv}[1]{\mathbf{D}^{(#1)}}
\newcommand{\tenp}[2]{\underset{#1,#2}{*}}
\title{Approximations of the Mortensen observer using higher order extended Kalman filters}
\author{Tobias Breiten${}^{\star}$ \and Justus Ramme${}^\star$ \and Jesper Schr\"oder${}^\dagger$}
\date{Source: arXiv:2604.25643v1 (CC BY 4.0)}
\begin{document}
\maketitle

\noindent\emph{Source and licence.} Approximations of the Mortensen observer using higher order extended Kalman filters. Version arXiv:2604.25643v1 (CC BY 4.0), distributed by arXiv under the Creative Commons Attribution 4.0 International licence, http://creativecommons.org/licenses/by/4.0/, as recorded in the arXiv licence metadata for this exact version and shown on the abstract page https://arxiv.org/abs/2604.25643v1. Full licence notice: \texttt{licenses/arxiv-2604.25643-CC-BY-4.0.txt}.\par\smallskip

\noindent\textit{Editorial scope.} Counted unit: the article's lemma LEM:HigherDerivativeHJB (source0.tex lines 656--671). The article's complete original proof of it is delivered with it (source0.tex lines 672--697, one \texttt{proof} environment of 26 lines). No mathematics is written, repaired or re-derived: the statement and the proof are the article's own lines, in the article's own notation, and no retained line is substituted or re-flowed.  Retained uncounted as the source-local context the proof needs: lines 247--255; lines 481--486; lines 531--541.  Nothing was omitted from any retained span, so the package carries no omission record.  No retained line contains a citation, so the wrapper carries no bibliography and no bibliography entry.  No retained line contains a figure environment, an image-inclusion command or drawing code, so no image is delivered and the package declares no asset dependency.  Editorial formatting only, in addition: an AMS \texttt{amsart} preamble that restates the article's own declarations and macros used by the retained lines, namely the environments \texttt{lemma}, \texttt{assumption} on separate counters, together with the standard packages those lines need.  The retained environments print in this document as Lemma 1, Assumption 1 in order of appearance, whereas the article prints the same items with the numbering of its own full document.  The complete native source of the article as distributed in the arXiv e-print for this exact version is bundled unchanged as \texttt{sources/199316/source0.tex}.
\begin{equation}\label{eq: HJB}
	\begin{aligned}
		\partial_t \mathcal{V}(t,\xi) &=
		-\nabla \mathcal{V}(t,\xi)^\top f(\xi) 
		- \frac{1}{2} \Vert G^\top \nabla \mathcal{V}(t,\xi) \Vert_{R^{-1}}^2
		+ \frac{1}{2} \Vert y(t) - h(\xi) \Vert_Q^2,\\
		\mathcal{V}(0,\xi) &= \frac{1}{2} \Vert \xi - x_0 \Vert_\Gamma^2,
	\end{aligned}
\end{equation}

\begin{lemma}\label{LEM:ProductRule}
Let $u,g\in C^k(\R^n;\R^m)$, then it holds that
  \begin{equation}\label{EQ:ProductRule}
    \TenDeriv{k}(u^\top g) = \sum_{i=0}^{k} \Sym_{i,k-i}\left( \TenDeriv{i}u \tenp{1}{1} \TenDeriv{k-i}g \right).
  \end{equation}
\end{lemma}

\begin{assumption}~ \label{ass:smoothness}
    \begin{enumerate}
        \item The functions $f \colon \mathbb{R}^n \to \mathbb{R}^n$ and $h \colon \mathbb{R}^n \to \mathbb{R}^p$ are $k$-times continuously differentiable.
        %
        \item For all $t \in [0,T]$ the value function $\mathcal{V}(t,\cdot)$ is $(k+1)$-times continuously differentiable.
        %
        \item The value function and its spatial derivatives are $C^1$ in time and the order of derivatives can be exchanged, i.e., for $\ell \leq k $ it holds
        $ \partial_t \TenDeriv{\ell} \mathcal{V}(\cdot,\widehat{x}_\mathrm{M}(\cdot)) = \TenDeriv{\ell} \partial_t \mathcal{V}(\cdot,\widehat{x}_\mathrm{M}(\cdot)) $.
    \end{enumerate}
    %
\end{assumption}

\begin{lemma}\label{LEM:HigherDerivativeHJB}
    Let \Cref{ass:smoothness} hold. Then for all $3 \leq j \leq k$ it holds
    %
    \begin{equation}
    \begin{aligned}\label{EQ:HigherDerivativeHJB}
        \partial_t \TenDeriv{j} \mathcal{V}(t,\xi) &= -\sum_{i=0}^{j} \Sym_{i,j-i}\left( \TenDeriv{i+1} \mathcal{V}(t,\xi) \tenp{1}{1}\TenDeriv{j-i}f(\xi) \right)
        \\ 
        &\qquad - \frac{1}{2}\sum_{i=0}^{j} \Sym_{i,j-i}\left( \TenDeriv{i+1} \mathcal{V}(t,\xi) \tenp{1}{1} \left( G R^{-1}G^\top \tenp{2}{1}\TenDeriv{j-i+1} \mathcal{V}(t,\xi) \right) \right)
        \\
        &\qquad - \TenDeriv{j} h(\xi) \tenp{1}{1} Q y(t) 
        \\
        & \qquad + \frac{1}{2} \sum_{i=0}^{j} \Sym_{i,j-i} \left(\TenDeriv{i} h(\xi) \tenp{1}{1} \left(Q \tenp{2}{1} \TenDeriv{j-i} h(\xi)\right) \right).
    \end{aligned}
    \end{equation}
    %
\end{lemma}

\begin{proof}
  We differentiate the three summands of \eqref{eq: HJB} separately using \Cref{LEM:ProductRule}. 
  For the first term, we obtain
  \begin{align*}
    \TenDeriv{j} \left(  \nabla \mathcal{V}(t,\xi)^\top f(\xi)\right) =& \sum_{i=0}^{j} \Sym_{i,j-i} \left( \TenDeriv{i} \left( \nabla \mathcal{V}(t,\xi) \right) \tenp{1}{1} \TenDeriv{j-i}f(\xi) \right) \\
    =&  \sum_{i=0}^{j} \Sym_{i,j-i} \left( \TenDeriv{i+1} \mathcal{V}(t,\xi) \tenp{1}{1} \TenDeriv{j-i}f(\xi) \right).
  \end{align*}
  Similarly, for the second term, we see that
  \begin{align*}
    \TenDeriv{j}\bigl(\Vert G^\top \nabla \mathcal{V}(t,\xi) \Vert_{R^{-1}}^2 \bigr) &=
    \TenDeriv{j}\bigl( \nabla \mathcal{V}(t,\xi)^\top  (  G R^{-1}G^\top  \nabla \mathcal{V}(t,\xi) ) \bigr) \\
    &= \sum_{i=0}^{j} \Sym_{i,j-i} \left( \TenDeriv{i} \left( \nabla \mathcal{V}(t,\xi) \right) \tenp{1}{1} \TenDeriv{j-i}\left( G R^{-1}G^\top  \nabla \mathcal{V}(t,\xi) \right)  \right)\\
    &= \sum_{i=0}^{j}\Sym_{i,j-i} \left( \TenDeriv{i+1}  \mathcal{V}(t,\xi) \tenp{1}{1} \left(G R^{-1}G^\top \tenp{2}{1}\TenDeriv{j-i+1} \mathcal{V}(t,\xi) \right)  \right).
  \end{align*}
  Finally differentiating the third term, we obtain
  \begin{align*}
      \TenDeriv{j}\left(\Vert y(t) - h(\xi) \Vert_{Q}^2\right) 
      &= \TenDeriv{j}\left( y(t)^T Q y(t) -2  h(\xi)^\top Q y(t) + h(\xi)^\top Q h(\xi) \right)
      \\
      &= \TenDeriv{j}\left( y(t)^T Q y(t)\right) -2  \TenDeriv{j}\left(h(\xi)^\top Q y(t)\right) + \TenDeriv{j}\left(h(\xi)^\top Q h(\xi) \right)
      \\ 
      &= -2   \TenDeriv{j} h(\xi) \tenp{1}{1} Q y(t) 
      \\
      & \qquad + \sum_{i=0}^{j} \Sym_{i,j-i} \left(\TenDeriv{i} h(\xi) \tenp{1}{1} \left(Q \tenp{2}{1} \TenDeriv{k-i} h(\xi)\right) \right).
  \end{align*}
\end{proof}

\end{document}
